% ============================================================
%  第2章：相关理论与技术基础
% ============================================================
\chapter{相关理论与技术基础}

\section{自注意力机制}

\zhlipsum[1]

Vaswani等人\cite{vaswani2017attention}提出注意力机制的核心计算过程，可表示为对查询（Query）、键（Key）、值（Value）三组向量的映射操作：
\begin{equation}
  \text{Attention}(Q, K, V) = \text{softmax}\!\left(\frac{QK^{T}}{\sqrt{d_k}}\right) V
  \label{eq:attention}
\end{equation}
其中$d_k$为键向量的维度，除以$\sqrt{d_k}$起到缩放作用。

\zhlipsum[2]

多头注意力机制将输入投影到$h$个不同的子空间，在每个子空间中独立计算注意力，最后拼接输出：
\begin{align}
  \text{MultiHead}(Q, K, V) &= \text{Concat}(\text{head}_1, \ldots, \text{head}_h) W^{O} \\
  \text{head}_i &= \text{Attention}(QW_i^{Q}, KW_i^{K}, VW_i^{V})
\end{align}

\section{Transformer架构}

\zhlipsum[14]
Transformer 作为一种完全基于注意力机制的序列转换模型，在多个自然语言处理任务上取得了突破性成果\cite{vaswani2017attention}。

每个子层均采用残差连接和层归一化：
\begin{equation}
  \text{LayerNorm}(x + \text{Sublayer}(x))
  \label{eq:layernorm}
\end{equation}

前馈网络由两个线性变换和一个非线性激活函数组成：
\begin{equation}
  \text{FFN}(x) = \max(0, xW_1 + b_1)W_2 + b_2
  \label{eq:ffn}
\end{equation}

\zhlipsum[5]

\section{预训练与微调范式}

\zhlipsum[7]
BERT 通过掩码语言模型与下一句预测任务进行预训练，在十一个自然语言处理任务上刷新了最优结果\cite{devlin2019bert}。

\section{高效注意力机制}

\zhlipsum[9]

\begin{equation}
  C_{\text{full}} = O(n^2 \cdot d)
  \label{eq:complexity_full}
\end{equation}

线性化注意力的目标是将复杂度降至$O(n \cdot d)$，代表性的方法包括稀疏注意力、低秩近似和核方法等\cite{vaswani2017attention}。
